% =============================================================================
% Chapter 9 --- Conclusion
% =============================================================================

\chapter{Conclusion}
\label{ch:conclusion}

This report has presented the RTL design and functional verification of
a controller for an in-memory-computing Artificial Neural Network
accelerator, carried out as a final project in the Technion VLSI
Laboratory. The design comprises three synthesizable modules ---
a combinational host interface, a hierarchical FSM controller with
three cooperating sub-FSMs for programming, verification, and erase,
and a synchronous input buffer --- arranged around a single
transaction-word format and a multi-bit pulse-mode bus that drives a
Y-Flash floating-gate memristor array.

The architectural contribution of the work is a closed-loop
program--verify--(re-program or erase) recovery policy implemented as
an FSM hierarchy: a single host-issued programming command is
translated by the controller into an iterative sequence of analog
pulses and verify reads, with re-programming on under-programmed
readbacks within a configurable retry budget and a fall-back to erase
on over-programmed readbacks or after retry exhaustion. The
controller absorbs device-characterization data through a Look-Up
Table that shapes the first-attempt programming pulse on a per-weight
basis, with a short single-cycle nudge train reserved for subsequent
corrections.

The methodological contribution of the work is a verification strategy
that establishes the digital correctness of this recovery policy
without access to the physical analog core. A behavioral mock of the
ANN array exposes the same digital boundary contract that the real
core will eventually present --- a shadow weight matrix, an
analog-side readback line, and an \emph{op\_done} handshake --- and
deliberate fault injection during the verify read window forces the
controller into every branch of the recovery diagram. The full
regression converges every weight in the matrix to its expected value
in the presence of injected failures, while the targeted directed
tests confirm the static transaction-word contract, the dynamic FSM
protocol of the three-module stack, the spatial selectivity of the
erase path, and the bit-serial inference dataflow. Together, these
results establish that the controller's digital behavior conforms to
its architectural intent across every claim of the verification plan.

The boundaries of the work are equally explicit. Behavioral mocks
stand in for the physical memristor array; \emph{op\_done} and the
readback line are idealized as deterministic functions of the FSM
state; a handful of architectural hooks in the RTL (the
currently-unreachable reset state, the two unused module parameters,
the latent status signals) are deliberately left as integration
points for a follow-on revision; and no synthesis, timing closure, or
power/area artifacts are produced. A rigorous next phase is therefore
not simply ``more tests'', but a graduated extension of the
verification: SystemVerilog assertions on the protocol boundaries,
constrained-random stimulus with quantitative coverage of the
recovery branches, and --- when the device team can supply it --- a
higher-fidelity model of the Y-Flash cell that replaces the
behavioral mock without disturbing the surrounding regression
infrastructure~\cite{Ielmini2018RRAM, Sebastian2020IMC,
Wang2022YFlash}.

The composability of the present environment is what makes this
progression realistic. The mock-based regression has been designed so
that every directed test that passes today will continue to run
against a higher-fidelity model tomorrow, and the recovery-branch
coverage established today carries forward into any future
mixed-signal phase. The controller delivered here is therefore not
merely a snapshot of digital behavior at one point in time; it is the
defensible digital baseline against which the analog integration of
the project's broader Y-Flash research effort can be carried out.
